\documentclass[a4paper, 12pt]{article}
% math symbols
\usepackage{amssymb}
\usepackage{amsmath}
\usepackage{mathrsfs}
\usepackage{mathseries}


\usepackage[margin = 2cm]{geometry}

\tolerance = 1000
\emergencystretch = 0.74cm



\pagestyle{empty}
\parindent = 0mm

\renewcommand{\coursetitle}{DM/ML}
\setcounter{curtask}{1}

\setmathstyle{\includegraphics[scale = 0.05]{../pics/utia-rest.png}}{Мат. логика}{3 курс}


\begin{document}

\libproblem{math-logic}{ef-r-q}
\libproblem{math-logic}{ef-n-n-plus-z}

\libproblem{math-logic}{ef-infinite-counterex}

\dzcomment{
    В одном из примеров достаточно очень маленького числа шагов игры, а значит и в сигнатуре имеет смысл
    рассматривать только предикаты малой арности (в одном из примеров).
}

\libproblem{math-logic}{ef-finite-isomorphism}
\libproblem{math-logic}{ef-orders-even}
\libproblem{math-logic}{ef-connected}



\libdefinition{theory-categorical}

\libproblem{math-logic}{bijection-without-cycles}
\libproblem{math-logic}{group-categorical}


\end{document}



%%% Local Variables:
%%% mode: latex
%%% TeX-master: t
%%% End:
